\clearpage
\appendix

\section{中心胞元表示定理的证明}
\label{app:center-cell-representation}

\begin{proof}[定理~\ref{thm:center-cell-representation} 的证明]
  在每个 $\partial\Omega_{0,\ell}$ 上，以下用上标 $+$ 表示从介质柱外部取得的迹，
  用上标 $-$ 表示从介质柱内部取得的迹；这两个上标与左右人工边界
  $\Gamma_-$、$\Gamma_+$ 无关。法向 $\nu$ 始终由介质柱内部指向外部。

  先确定均匀背景场。由于介质柱与 $\Gamma_\pm$ 严格分离，$u$ 在两个人工边界的
  邻域内满足波数为 $\kappa_{\mathrm e}$ 的均匀 Helmholtz 方程。将 $u$ 及
  $\partial_xu$ 的迹投影到 $\psi_m$ 后，每个 Fourier 分量的两个 Rayleigh 系数由
  \begin{equation*}
    \begin{bmatrix}
      1&1\\ \ii\gamma_m&-\ii\gamma_m
    \end{bmatrix}
  \end{equation*}
  所给出的二元关系确定。因为 $\gamma_m\neq0$，该矩阵可逆，所以锚定在
  $\Gamma_-$ 的右行系数 $a_{0,m}^-$ 与锚定在 $\Gamma_+$ 的左行系数
  $b_{0,m}^+$ 均由 $u$ 唯一确定。用这些系数构造正文中的
  $u_{\mathrm{hom}}$，并在中心胞元的背景区域令
  \begin{equation*}
    w^+:=u-u_{\mathrm{hom}}.
  \end{equation*}
  按构造，$w^+$ 在左端没有右行入射分量，在右端没有左行入射分量，因而在两个端口都满足
  outgoing Rayleigh 条件。对第 $\ell$ 根介质柱，再令 $w_\ell^-:=u|_{\Omega_{0,\ell}}$。

  接着构造层势密度。对每个 $\ell$，考虑交换内外波数的辅助问题：在
  $\Omega_{0,\ell}$ 内使用 $\kappa_{\mathrm e}$，在
  $\RR^2\setminus\overline{\Omega_{0,\ell}}$ 内使用
  $\kappa_{\mathrm i,\ell}$，并要求外部场满足 outgoing Sommerfeld 条件。其界面跳跃取为
  \begin{align*}
    v_\ell^+-v_\ell^-
    &=-w^+-w_\ell^-,\\
    \partial_\nu v_\ell^+-\partial_\nu v_\ell^-
    &=-\partial_\nu w^+-\partial_\nu w_\ell^-.
  \end{align*}
  定理的非例外性假设保证该问题具有唯一解 $v_\ell$。定义
  \begin{align*}
    \tau_\ell
    &:=w^+-v_\ell^-
      =-w_\ell^--v_\ell^+,\\
    \sigma_\ell
    &:=-\partial_\nu w^++\partial_\nu v_\ell^-
      =\partial_\nu w_\ell^-+\partial_\nu v_\ell^+.
  \end{align*}
  两组等式分别由辅助问题的两个跳跃条件得到。

  现在验证这些密度给出所需场。自由空间 Green 表示公式及其消去公式表明，对
  $\mathbf z\in\Omega_{0,\ell}$，
  \begin{equation*}
    \left(
      \mathcal S_\ell^{(\kappa_{\mathrm i,\ell})}[\sigma_\ell]
      +\mathcal D_\ell^{(\kappa_{\mathrm i,\ell})}[\tau_\ell]
    \right)(\mathbf z)
    =w_\ell^-(\mathbf z).
  \end{equation*}
  事实上，代入 $\tau_\ell$、$\sigma_\ell$ 的第二组表达式后，来自
  $w_\ell^-$ 的两项由内部 Green 公式重构 $w_\ell^-$，来自外部场
  $v_\ell^+$ 的两项则在介质柱内部相消。

  对背景区域，在
  $\mathcal C_0\setminus\overline{\Omega_0}$ 上对 $w^+$ 与
  $G_{\mathrm{QP}}^{(\kappa_{\mathrm e})}$ 使用 Green 第二恒等式。顶、底边界积分因
  $\beta$-准周期相位互相抵消；两个竖直端口上的积分在逐项代入 outgoing Rayleigh 展开后
  也为零。这里每个 Fourier 分量的消去因子是同一个 $\gamma_m$，而
  $\gamma_m\neq0$ 保证系数分解没有退化。因此
  \begin{equation*}
    \sum_{\ell=1}^{N_0}\left(
      -\mathcal S_{\mathrm{QP},\ell}^{(\kappa_{\mathrm e})}
        [\partial_\nu w^+]
      +\mathcal D_{\mathrm{QP},\ell}^{(\kappa_{\mathrm e})}[w^+]
    \right)
    =w^+.
  \end{equation*}
  对 $v_\ell^-$ 使用内部 Green 公式及其在介质柱外部的消去公式，则有
  \begin{equation*}
    \mathcal S_{\mathrm{QP},\ell}^{(\kappa_{\mathrm e})}
      [\partial_\nu v_\ell^-]
    -\mathcal D_{\mathrm{QP},\ell}^{(\kappa_{\mathrm e})}[v_\ell^-]
    =0
    \qquad\text{在 }\mathcal C_0\setminus\overline{\Omega_0}.
  \end{equation*}
  将 $\tau_\ell$、$\sigma_\ell$ 的第一组表达式代入并对 $\ell$ 求和，便得到
  \begin{equation*}
    \sum_{\ell=1}^{N_0}\left(
      \mathcal S_{\mathrm{QP},\ell}^{(\kappa_{\mathrm e})}[\sigma_\ell]
      +\mathcal D_{\mathrm{QP},\ell}^{(\kappa_{\mathrm e})}[\tau_\ell]
    \right)
    =w^+.
  \end{equation*}
  再加回 $u_{\mathrm{hom}}$，即得定理所述的外部表示；结合前面的内部表示，存在性得证。

  最后讨论唯一性。两组表示之差仍具有定理中的形式，并表示零场。前述端口系数提取关系立即
  给出所有 $a_{0,m}^-$、$b_{0,m}^+$ 均为零，所以 Rayleigh 系数列无条件唯一。若再满足定理
  末尾的互补系统零核条件，则只剩的层势部分在交换适用区域并改变内部层势符号后，由层势
  跳跃关系给出该齐次互补系统的一个解；它在每个界面上的共同 Dirichlet 和 Neumann 迹分别
  为 $-\tau_\ell$ 和 $\sigma_\ell$。零核条件迫使互补场恒为零，再次使用层势跳跃关系便有
  $\tau_\ell=0$、$\sigma_\ell=0$，$1\leq\ell\leq N_0$。因此在这一附加条件下，全部层势
  密度也唯一。
\end{proof}
